\documentclass[11pt, a4paper]{report}

% ─── Encoding & language ────────────────────────────────────────────────────
\usepackage[utf8]{inputenc}
\usepackage[T1]{fontenc}
\usepackage[dutch]{babel}

% ─── Layout & styling ───────────────────────────────────────────────────────
\usepackage{geometry}
\usepackage{xcolor}
\usepackage{graphicx}
\usepackage{fancyhdr}
\usepackage{titlesec}
\usepackage{tocloft}
\usepackage{parskip}
\usepackage{emptypage}

% ─── Tables & lists ─────────────────────────────────────────────────────────
\usepackage{booktabs}
\usepackage{longtable}
\usepackage{enumitem}
\usepackage{array}

% ─── Hyperlinks ─────────────────────────────────────────────────────────────
\usepackage{hyperref}

% ═══════════════════════════════════════════════════════════════════════════
%  COLOUR PALETTE  (Avans Hogeschool rood)
% ═══════════════════════════════════════════════════════════════════════════
\definecolor{avansred}{RGB}{193, 0, 0}
\definecolor{avansdark}{RGB}{130, 0, 0}
\definecolor{avanslight}{RGB}{245, 215, 215}
\definecolor{textgray}{RGB}{80, 80, 80}
\definecolor{rulegray}{RGB}{200, 200, 200}

% ─── Hyperref setup ─────────────────────────────────────────────────────────
\hypersetup{
    colorlinks  = true,
    linkcolor   = black,
    urlcolor    = avansdark,
    citecolor   = avansdark,
    pdftitle    = {Onderhoudbaarheid Rapport},
    pdfauthor   = {Avans Hogeschool},
}

% ─── Page geometry ──────────────────────────────────────────────────────────
\geometry{
    a4paper,
    top    = 2.5cm,
    bottom = 2.5cm,
    left   = 2.5cm,
    right  = 2.5cm,
}
\setlength{\headheight}{14.5pt}

% ─── Header / footer ────────────────────────────────────────────────────────
\pagestyle{fancy}
\fancyhf{}
\fancyhead[L]{\small\textbf{Avans Hogeschool}}
\fancyhead[R]{\small\color{textgray}\nouppercase{\leftmark}}
\fancyfoot[C]{\small\color{textgray}\thepage}
\renewcommand{\headrulewidth}{0.6pt}
\renewcommand{\headrule}{%
    \hbox to\headwidth{\color{avansred}\leaders\hrule height \headrulewidth\hfill}}
\renewcommand{\footrulewidth}{0pt}

% ─── Chapter titles ─────────────────────────────────────────────────────────
\titleformat{\chapter}[hang]
    {\normalfont\LARGE\bfseries}
    {\thechapter}
    {0.8em}
    {}
    [\vspace{6pt}\color{avansred}\rule{\linewidth}{1.2pt}]

\titlespacing*{\chapter}{0pt}{20pt}{12pt}

% ─── Section / subsection titles ────────────────────────────────────────────
\titleformat{\section}
    {\large\bfseries}
    {\thesection}{1em}{}

\titleformat{\subsection}
    {\normalsize\bfseries}
    {\thesubsection}{1em}{}

% ─── Table of contents styling ───────────────────────────────────────────────
\renewcommand{\cftchapfont}{\bfseries}
\renewcommand{\cftchappagefont}{\bfseries}
\renewcommand{\cftsecfont}{\color{textgray}}
\renewcommand{\cftsecpagefont}{\color{textgray}}
\setlength{\cftbeforechapskip}{6pt}

% ════════════════════════════════════════════════════════════════════════════
\begin{document}

% ─── Title page ─────────────────────────────────────────────────────────────
\begin{titlepage}
    \centering
    \vspace*{0.8cm}

    % Top rule
    {\color{avansred}\rule{\linewidth}{3pt}}\\[0.4cm]
    {\color{avansred}\rule{\linewidth}{0.8pt}}\\[1.4cm]

    % Document type label
    {\large\color{textgray}\textsc{Eindrapport --- Webservices REST Audit}}\\[0.8cm]

    % Main title
    {\fontsize{34}{40}\selectfont\bfseries\color{avansred} Onderhoudbaarheid}\\[0.4cm]
    {\Large\color{avansdark} Analyse, verbeteringen \& validatie}\\[1.8cm]

    % Meta table
    \begin{tabular}{@{}ll@{}}
        \color{textgray}\textbf{Auteurs}      & [Naam 1]                  \\[2pt]
                                               & [Naam 2]                  \\[2pt]
                                               & [Naam 3]                  \\[10pt]
        \color{textgray}\textbf{Datum}         & \today                    \\[2pt]
        \color{textgray}\textbf{Versie}        & 1.0                       \\[2pt]
        \color{textgray}\textbf{Opleiding}     & HBO-ICT                   \\[2pt]
        \color{textgray}\textbf{Instelling}    & Avans Hogeschool          \\[2pt]
        \color{textgray}\textbf{Leerjaar}      & Jaar 2 --- Periode 4      \\
    \end{tabular}

    \vfill

    % Bottom rule
    {\color{avansred}\rule{\linewidth}{0.8pt}}\\[0.3cm]
    {\color{avansred}\rule{\linewidth}{3pt}}\\[0.6cm]

    {\small\color{textgray}Avans Hogeschool, Breda}
\end{titlepage}

% ─── Table of contents ──────────────────────────────────────────────────────
\tableofcontents
\clearpage

% ════════════════════════════════════════════════════════════════════════════
%  CHAPTERS
% ════════════════════════════════════════════════════════════════════════════

% ─── 1 ──────────────────────────────────────────────────────────────────────
\chapter{Analyse onderhoudbaarheid}

\section{Inleiding}
% TODO: Beschrijf het doel van de analyse en wat verstaan wordt onder
%       onderhoudbaarheid (bijv. ISO 25010).

\section{Meetmethode}
% TODO: Beschrijf welke tools en methoden gebruikt zijn voor de analyse
%       (bijv. SonarQube, CodeClimate, handmatige beoordeling).

\section{Resultaten}
% TODO: Presenteer de resultaten van de analyse, inclusief metrieken.

\section{Conclusie}
% TODO: Sluit het hoofdstuk af met een oordeel over de huidige staat.


% ─── 2 ──────────────────────────────────────────────────────────────────────
\chapter{Testopzet en testresultaten}

\section{Inleiding}
% TODO: Beschrijf de doelstelling van de testfase.

\section{Testopzet}
% TODO: Beschrijf welke soorten tests zijn opgezet (unit, integratie,
%       end-to-end) en welke tooling gebruikt is.

\section{Testresultaten}
% TODO: Presenteer de resultaten per testsoort, inclusief coverage.

\section{Analyse van de resultaten}
% TODO: Interpreteer de resultaten en leg verbanden met de onderhoudbaarheidsscore.

\section{Conclusie}
% TODO: Sluit het hoofdstuk af.


% ─── 3 ──────────────────────────────────────────────────────────────────────
\chapter{Verbeteringen (prioritering en onderbouwing)}

\section{Inleiding}
% TODO: Beschrijf de aanpak voor het bepalen en prioriteren van verbeteringen.

\section{Overzicht van verbeteringen}
% TODO: Geef een overzicht van alle voorgestelde verbeteringen.

\section{Prioritering}
% TODO: Leg uit op welke basis de prioritering is bepaald
%       (bijv. impact vs. inspanning, MoSCoW).

\section{Onderbouwing}
% TODO: Onderbouw de keuzes per verbetering met data of argumentatie.

\section{Conclusie}
% TODO: Sluit het hoofdstuk af.


% ─── 4 ──────────────────────────────────────────────────────────────────────
\chapter{Aangepast ontwerp}

\section{Inleiding}
% TODO: Beschrijf waarom het ontwerp aangepast is en wat de uitgangspunten zijn.

\section{Huidig ontwerp}
% TODO: Beschrijf het huidige ontwerp (architectuur, structuur, patronen).

\section{Nieuw ontwerp}
% TODO: Beschrijf het aangepaste ontwerp en motiveer de keuzes.

\section{Vergelijking}
% TODO: Vergelijk het oud en nieuw ontwerp op meetbare punten.

\section{Conclusie}
% TODO: Sluit het hoofdstuk af.


% ─── 5 ──────────────────────────────────────────────────────────────────────
\chapter{Realisatie (PoC) \& verantwoording}

\section{Inleiding}
% TODO: Beschrijf het doel van de Proof of Concept.

\section{Scope van de PoC}
% TODO: Beschrijf welke verbeteringen in de PoC zijn meegenomen en welke niet.

\section{Implementatie}
% TODO: Beschrijf de gemaakte wijzigingen en hoe deze zijn doorgevoerd.

\section{Verantwoording van keuzes}
% TODO: Onderbouw de technische keuzes die tijdens de realisatie zijn gemaakt.

\section{Conclusie}
% TODO: Sluit het hoofdstuk af met een beoordeling van de PoC.


% ─── 6 ──────────────────────────────────────────────────────────────────────
\chapter{Validatie verbeteringen (testen \& regressie)}

\section{Inleiding}
% TODO: Beschrijf het doel van de validatiefase.

\section{Validatieaanpak}
% TODO: Beschrijf hoe de verbeteringen gevalideerd zijn.

\section{Testresultaten na verbetering}
% TODO: Presenteer de hertest-resultaten en vergelijk met de beginsituatie.

\section{Regressietests}
% TODO: Beschrijf welke regressietests uitgevoerd zijn en wat de uitkomst is.

\section{Eindconclusie}
% TODO: Sluit het rapport af met een eindoordeel over de verbetering van
%       de onderhoudbaarheid.


% ════════════════════════════════════════════════════════════════════════════
\end{document}
